\subsection{Ideal complete-view privacy and modular cryptographic composition}\label{sec:zkrs:composition}
The public core is not the complete verifier view.
A real implementation includes commitments to packed witness elements,
hidden pads, correction coefficients and auxiliary masks, plus the
affine-proof transcripts and any authenticated oracle queries.
We now state exactly what an external proof system must supply.

\begin{definition}[Opaque commitment-and-affine ideal functionality]\label{def:zkrs:ideal-func}
An \emph{opaque ideal proof interface} stores hidden witness data in a
private registry, but its public transcript has only the following effects:
\begin{enumerate}[label=\textup{(\roman*)}]
\item On binding a hidden tuple, it emits a public registration token
sampled from a specified distribution independent of the hidden tuple.
The token is produced before the challenges that rely on its binding.
\item Given a proposed public affine relation, it privately evaluates
the relation against the registered hidden tuple and releases its
accept/reject bit, with no other witness-dependent public data.
\item It never provides the verifier a raw hidden field value, an
unmasked codeword symbol, or an authenticated opening of a hidden value.
\item On honest admissible executions it always outputs acceptance.
For malicious executions it enforces the registered affine relations
exactly. Its public token values may be adaptively queried or reused
only to the extent explicitly supported by the ideal interface.
\end{enumerate}
This is a \emph{mathematical ideal functionality}, not an assertion
that a concrete ordinary binding commitment is perfectly opaque.
\end{definition}

\begin{lemma}[Opaque tokens do not perturb the public-core law]\label{lem:zkrs:opaque-kernels}
Assume all opaque interface tokens and outputs for an
honest execution are independent of the private witness and
of the random masks used in the public core, except that the
acceptance bits are the constant $1$.
Then, conditioned on any allowed visible history of tokens,
the honest public-core transitions have the same uniform
conditional laws as Lemmas~\ref{lem:zkrs:input-uniform},
\ref{lem:zkrs:U-conditional}, and~\ref{lem:zkrs:A-conditional}.
\end{lemma}
\begin{proof}
Independence implies that conditioning on the tokens
does not change the joint law of the private masks
and the verifier's previous public-core messages.
Constant acceptance bits reveal no new information.
The input-pair translation argument, trace-zero
pushforward argument, and final uniform-$\rho$
translation argument therefore hold under the
additional conditioning without change.
These are exactly the three cited lemmas.
\end{proof}

\begin{definition}[Complete ideal verifier view]\label{def:zkrs:ideal-view}
For a single admissible opening, the \emph{complete ideal view}
consists of the public commitment tokens, the verifier's own
random tape, $(m,y')$, its challenges $(\theta,z)$,
the public vector $U$, the final scalar $A$,
the public trace checks, the acceptance bits of all
required affine relations, and all interaction with
the opaque ideal interface.
The verifier may choose its challenges adaptively
after all messages visible at that stage.
\end{definition}

\begin{theorem}[Perfect complete-view ZK in the opaque ideal model]\label{thm:zkrs:ideal-full-zk}
For any causal verifier and any two admissible witnesses
for the same fixed public evaluation instance, the complete
ideal verifier views of Definition~\ref{def:zkrs:ideal-view}
are identically distributed. A simulator without the witness
generates the opaque registration tokens from their specified
distributions, uses Definition~\ref{def:zkrs:core-sim}
for the public core, and returns the constant honest
acceptance bits from the ideal checks.
\end{theorem}
\begin{proof}
The simulator samples each token from the
witness-independent distribution prescribed by
Definition~\ref{def:zkrs:ideal-func}, and couples
the verifier's random tape with the real experiment.
By Lemma~\ref{lem:zkrs:opaque-kernels},
conditioning on these tokens preserves the
uniform laws of all public core messages.
Theorem~\ref{thm:zkrs:perfect-core}
then gives identical conditional distributions
at every public-core round, including when
the verifier chooses challenges from the
earlier tokens and messages.
The honest hidden state satisfies all
ideal affine relations by
Lemmas~\ref{lem:zkrs:honest-checker} and
\ref{lem:zkrs:honest-relations};
hence their public accept/reject bits are
the constant $1$ in both experiments.
All outputs are either the same coupled
verifier actions, identically distributed
tokens, identically distributed core messages,
or these deterministic accept bits.
Thus the complete finite joint view laws
coincide.
\end{proof}

\begin{definition}[Sequential external-interface simulation contract]\label{def:zkrs:realization}
Fix a finite interactive implementation, a verifier class and
a bound on authenticated queries. Resolve each protocol step into
a finite transition on its \emph{visible history}, which records the
entire public prefix and the verifier's random tape.
For a fixed admissible witness $w$, let
$K_i^w(h)$ be the real conditional distribution
of the next visible state from history $h$.
Let $L_i(h)$ be a causal, \emph{witness-independent}
transition kernel producing an ideal opaque-interface
view as in Definition~\ref{def:zkrs:ideal-view}.
An \emph{online sequential realization with local errors}
$(\varepsilon_0,\ldots,\varepsilon_N)$ means:
\begin{enumerate}[label=\textup{(\roman*)}]
\item initial real and ideal visible-state distributions
have statistical distance at most $\varepsilon_0$;
\item for each stage $1\leq i\leq N$ and each possible
visible history $h$,
\[
 \Delta(K_i^w(h),L_i(h))\leq\varepsilon_i,
\]
uniformly over all admissible witnesses and
all verifiers in the chosen class;
\item the ideal kernels $L_i$ faithfully implement
the public-core simulator and opaque-interface
simulator, including the same ordering and
causal challenge dependence as the real protocol.
\end{enumerate}
This requires checking correlations with any
preceding commitments and authenticated openings:
the perfect public-core law of
Theorem~\ref{thm:zkrs:perfect-core}
does \emph{not} alone establish (ii)
for histories containing real commitment roots.
For histories that are unreachable under a
particular process, one may choose a regular
conditional kernel arbitrarily, but the uniform
bound must still hold under the chosen extension.
\end{definition}

\begin{theorem}[Conditional real-world statistical ZK composition]\label{thm:zkrs:conditional-zk}
Suppose a concrete single-opening implementation satisfies
Definition~\ref{def:zkrs:realization} with local errors
$(\varepsilon_0,\ldots,\varepsilon_N)$
for a given query bound and verifier class.
Then for every admissible witness,
the implemented complete view has a witness-independent
simulator with statistical distance at most
\[
 \boxed{
 \varepsilon_{\rm ZK}\leq
 \min\left\{1,\sum_{i=0}^{N}\varepsilon_i\right\}.
 }
\]
If all external interfaces are perfectly simulatable,
the complete protocol is perfect zero knowledge for
that single-opening game.
\end{theorem}
\begin{proof}
Use the finite visible history spaces of the
real and ideal processes, with the verifier's
random tape incorporated into each history.
Condition (ii) of
Definition~\ref{def:zkrs:realization}
bounds the total variation between the real
and ideal transition kernels at every stage
and every visible history.
Condition (i) bounds the initial-state
distance. The stepwise hybrid lemma,
Lemma~\ref{lem:security:hybrid},
therefore bounds the distance between the
final visible-history distributions by
$\sum_{i=0}^{N}\varepsilon_i$.
Condition (iii) and
Theorem~\ref{thm:zkrs:ideal-full-zk}
provide a witness-independent implementation
of the ideal visible process.
A deterministic projection from final
histories to the actual verifier view
cannot increase total variation, by
Corollary~\ref{cor:security:projection}.
Finally, total variation is at most
one for any two distributions.
\end{proof}

\begin{definition}[External algebraic-soundness failure]\label{def:zkrs:sound-failure}
Let $\mathcal B$ be the event that an
implemented verifier accepts despite
some underlying commitment-binding,
packed-witness consistency, or advertised
field-affine constraint failing to hold
for the fixed hidden tuple.
Assume a separately established
implementation-specific error bound
$\Pr[\mathcal B]\leq\varepsilon_{\rm ext,snd}$.
\end{definition}

\begin{theorem}[Conditional full implementation soundness]\label{thm:zkrs:conditional-sound}
Under the causal and fixed-before-challenge hypotheses
of Theorem~\ref{thm:zkrs:false-sound},
and assuming implementation acceptance
outside $\mathcal B$ implies acceptance of all ideal
algebraic checks on the same fixed hidden values, a false
evaluation claim is accepted with probability at most
\[
 \boxed{
 \varepsilon_{\rm snd}\leq
 \min\left\{1,\frac{\ell+n}{q}
             +\varepsilon_{\rm ext,snd}\right\}.
 }
\]
\end{theorem}
\begin{proof}
Outside $\mathcal B$, implementation
acceptance entails acceptance of the
ideal algebraic checker.
Thus the implementation-accept event
is contained in the union of the
ideal-accept event and $\mathcal B$.
Theorem~\ref{thm:zkrs:false-sound}
bounds the first event by
$(\ell+n)/q$, and
Definition~\ref{def:zkrs:sound-failure}
bounds the second by
$\varepsilon_{\rm ext,snd}$.
Apply the probability union bound
and cap the result by one.
\end{proof}

\begin{remark}[No unconditional real-world ZK claim]
Theorem~\ref{thm:zkrs:conditional-zk}
is deliberately conditional:
the implementation must actually
satisfy a simulation contract for
its \emph{full} commitment and proof
transcripts. Neither ordinary
computational hiding of one commitment
nor a proof that $(m,y',U,A)$ has
uniform marginals establishes that
contract. A real online commitment
chosen before the evaluation point
also requires a simulator capable of
maintaining consistency with later
chosen evaluations.
\end{remark}
